\scanpage{132}
U N I V E R S I T E\\
DES SCIENCES ET TECHNIQUES\\
DU LANGUEDOC\\
\texttt{------------}

Montpellier, le 25 Novembre 1981

Service de la Scolarité\\
MA / GJ\\
Diffusion : Liste A\\
N° : 6455

AVIS DE SOUTENANCE D'UNE THESE DE DOCTORAT DE 3ème CYCLE

CHIMIE ORGANIQUE PHYSIQUE\\
\texttt{-------------------------}

\uline{Nom et prénom du candidat} : LORENTE Philippe

\uline{Titre de la thèse} : HYDROLYSE D'ORTHOESTERS CYCLIQUES AROMATIQUES\\
Compétition entre effet stéréoélectronique et effet de groupe partant.

\uline{Date, heure et lieu de soutenance prévus} :\\
11 DECEMBRE 1981 à 16 Heures\\
Salle de Cours n° 2 - USTL

\uline{Nom du Directeur de thèse} :\\
Monsieur le Professeur G. LAMATY

\uline{Désignation du laboratoire où la thèse est préparée} :\\
Laboratoire de Chimie organique physique.

\uline{Résumé de la thèse} :

Quatre orthoesters cycliques aromatiques dérivés de la dihydro-3,4 coumarine substituée sur le noyau benzénique ont été synthétisés ; la réaction d'hydro-\\
lyse de ces composés a été étudiée dans le but d'opposer un effet de groupe partant à l'effet stéréoélectronique qui se manifeste dans ce type de réaction.Une évaluation de l'énergie due à l'effet stéréoélectronique pouvait résulter de l'étude de cette compé-\\
tition.

Deux hypothèses ont été formulées pour rendre compte des résultats expérimentaux ; l'une d'entre elle conduit à une estimation de l'énergie stéréoélectro-\\
nique.

\uline{Durée approximative des recherches ayant abouti à la thèse} : 2 ans

\uline{Composition du jury proposé} :\\
MM. : LAMATY Gérard\\
GENESTE Patrick\\
MOREAU Claude\\
VIALLEFONT Philippe

P. le Secrétaire Général,\\
l'Attachée d'Administration Universitaire,

\figure{\begin{tikzpicture}[x=0.8cm,y=0.8cm,line width=0.6pt]
\draw (0.48,-0.6) .. controls (0.44,-0.04) and (0.55,0.9) .. (0.56,1.12)
 .. controls (0.45,1.3) and (0.45,0.85) .. (0.51,0.61)
 .. controls (0.26,0.98) and (0.32,0.58) .. (0.57,0.64)
 .. controls (0.77,0.72) and (0.8,0.65) .. (0.93,0.64);
\draw (0.52,0.58) .. controls (-0.09,0.22) and (-0.04,0.02) .. (0.42,0.1)
 .. controls (0.65,0.14) and (0.91,0.21) .. (0.9,0.22)
 -- (0.64,0.16) .. controls (1.11,0.11) and (1.87,0.13) .. (2.07,0.2)
 .. controls (2.22,0.24) and (2.24,0.53) .. (2.16,0.68)
 .. controls (2.05,0.79) and (2.14,0.29) .. (2.17,0.01);
\draw (2.17,0.2) -- (2.17,-0.02);
\draw (3.02,-0.03) .. controls (2.8,-0.12) and (0.37,-0.11) .. (0.26,-0.24)
 .. controls (0.15,-0.33) and (0.5,-0.31) .. (0.55,-0.3);
\draw (-0.02,0.17)--(-0.08,0.23);
\end{tikzpicture}}
A. GENTIL

\mspage{131}
